% Results — figure-driven description of force-triggered stepping, descending
% arm (CUT->RG-E and BS->RG-{E,F} plastic).
%   fig_init_robustness      = round 6: 10 STDP initialisations, medium walk, lambda=1e-3
%   fig_stdp_weight_matrix   = (supplementary) weights, both legs x 3 projections x 4 modes
%   fig_stdp_weights_grid    = weights, 4 modes x 5 lambda (3 projections/panel)
%   fig_force_stages_lam*    = force at 3 learning stages, 4 modes, per lambda
%   fig_network_matrix_lam*  = 16 population rates x 4 modes, per lambda
%   fig_mode_lambda_*        = 4x5 summary heatmaps / trends / table
% Data: results/2026-09-25_ftgrid (run_forcetrig_grid.sh), round 6
% (results/2026-09-07 + results/2026-09-25, run_cutforce_sweep6.sh).

Before the individual results we fix the terms used throughout. The model is a
two-legged spinal CPG: a spiking network that, driven by a weak tonic brainstem
input and sensory feedback, produces the alternating extensor/flexor muscle
activity of walking. Each leg contains two mutually-inhibiting neuron pools ---
an \emph{extensor} and a \emph{flexor} rhythm generator (RG-E and RG-F) --- that
drive motor neurons and simple muscle models. Stance and swing are not imposed
by a clock: each leg's paw-contact (cutaneous) input is switched on when its own
extensor force rises and off when that force falls
(\autoref{sec:methods:paced}). A good walking rhythm is one in which extensor
and flexor are active in strict alternation (``counter-phase''); we measure this
by the Pearson correlation between the two muscle forces, which is $-1$ for
perfect alternation and $0$ for none, and measure left--right alternation the
same way between the two legs' extensor forces. Because each stance phase also
has a timeout, we check every condition for how often stance ends at that
timeout rather than on the force threshold (``frac.\ stance at cap''): near $0$
means the rhythm is genuinely produced by the closed sensorimotor loop, near $1$
means the timeout is acting as a clock. Three synapses are \emph{plastic} ---
the cutaneous CUT$\to$RG-E projection and the brainstem BS$\to$RG-E and
BS$\to$RG-F projections --- and change as the network runs, through STDP. The
Results ask how this self-tuning builds the gait, and how well the learned gait
holds up across walking speed and under partial unloading, as in
spinal-cord-injury rehabilitation.

\subsection{The learned set-point does not depend on the initial weights}\label{sec:results:robustness}

We first ask whether what the network learns depends on where it starts. We ran
medium walk at $\lambda=10^{-3}$ from ten initial weight distributions, from
$\mu=0$ (plastic synapses start silent) to $\mu=16$\,pA
(Figure~\ref{fig:init_robustness}). All ten runs converge to the same learned
weights within $\approx20$\,s: CUT$\to$RG-E to $60$--$62$\,pA and both
BS$\to$RG projections to $\approx16$\,pA, about half their $30$\,pA ceiling.
The gait they produce is equally independent of the starting point: at every
initialisation stance ends on the force threshold (no stance bout at the
timeout), extensor and flexor force alternate with corr$(F_E,F_F)$ between
$-0.67$ and $-0.70$ on both legs, and the two legs alternate with
corr$(F_{E,L},F_{E,R})$ between $-0.77$ and $-0.79$.

\begin{figure}[t]
  \centering
  \includegraphics[width=\textwidth]{fig_init_robustness.png}
  \caption{\textbf{Initialisation robustness of force-triggered stepping.}
    Medium walk, $\lambda=10^{-3}$, ten STDP initial-weight distributions
    ($\mu=0$--$16$\,pA; colour). \textbf{(a--c)} Weight trajectories of
    CUT$\to$RG-E, BS$\to$RG-E and BS$\to$RG-F (leg L, pA) versus time (s).
    \textbf{(d)} Steady-state ($t\ge30$\,s) gait metrics versus initial mean
    weight: extensor/flexor force correlation per leg, left/right extensor
    force correlation, and the fraction of stance bouts ended by the timeout.}
  \label{fig:init_robustness}
\end{figure}

\subsection{The learned synapses settle at values that do not depend on the walking mode}\label{sec:results:weights}

We next vary the locomotion mode and the learning rate together, on a grid of
four modes by five rates. The modes are slow, medium and fast walk (full
weight-bearing; medium walk is also called \emph{plantar stepping} or
\emph{baseline}) and \emph{toe stepping} (partial body-weight support, sensory
feedback halved). Their operating points (Table~\ref{tab:modes}) give strides of
$1016\pm36$, $812\pm33$ and $594\pm24$\,ms for slow, medium and fast walk and
$845\pm66$\,ms for toe stepping at $\lambda=10^{-3}$ (last 20\,s). The
differences come from stance ($417$, $312$, $188$ and $360$\,ms), since swing
always lasts its timeout (\autoref{sec:methods:paced}). The learning rate
$\lambda$ sets how much each plastic synapse changes per correlated pre/post
spike pair and is swept across five decades,
$\lambda\in\{10^{-2},10^{-3},10^{-4},10^{-5},10^{-6}\}$.

Once learning has converged, the synapses reach the same values in every mode
(Figure~\ref{fig:weights_grid}; both legs in Supplementary
Figure~\ref{fig:weight_matrix}): CUT$\to$RG-E settles at $61$--$63$\,pA and
BS$\to$RG-E/F at $15$--$16$\,pA, and the left and right legs learn the same
values. The learning rate sets how fast the synapses get there. At
$\lambda=10^{-3}$, CUT$\to$RG-E reaches $90\%$ of its final value within
$\approx6$\,s in the three walking modes; at $10^{-4}$ it takes $50$--$64$\,s. At
the fastest rate ($10^{-2}$) the synapses jump to their range within a few
seconds but then fluctuate from step to step below the plateau
(CUT$\to$RG-E $\approx46$--$58$\,pA). At $10^{-5}$ and $10^{-6}$ they barely
move within the $120$\,s run (CUT$\to$RG-E reaching $\approx7$--$14$\,pA and
staying at its $\approx4$\,pA initial value, respectively). Partial unloading
slows learning, because the paw-contact drive that trains CUT$\to$RG-E is itself
weaker: toe stepping needs $14$\,s instead of $6$\,s at $\lambda=10^{-3}$, and at
$10^{-4}$ it has reached only $\approx47$--$50$\,pA by the end of the run.

\begin{figure}[t]
  \centering
  \includegraphics[width=0.95\textwidth]{fig_stdp_weights_grid.png}
  \caption{\textbf{Weight trajectories of the three plastic projections, per
    locomotion mode and learning rate.} Rows: the four locomotion modes;
    columns: the five learning rates. Each panel shows CUT$\to$RG-E (solid,
    left axis, pA) and BS$\to$RG-E, BS$\to$RG-F (dashed, right axis, pA) for leg
    L versus time (s). Panels (a)--(t). The learning rate sets the time to
    plateau; partial unloading (bottom row) slows CUT$\to$RG-E convergence.}
  \label{fig:weights_grid}
\end{figure}

% =============================================================
\subsection{Rehabilitation progress: the gait forms as the weights converge}\label{sec:results:force}

We next examine whether the learned wiring produces good walking, and how the
gait builds up over a session. Figures~\ref{fig:force_l2}--\ref{fig:force_l6}
show the extensor and flexor muscle forces in three short windows --- the
beginning, middle and end of each $120$\,s run --- one figure per learning rate,
so that reading a figure left-to-right traces how the gait develops as the
synapses strengthen. With learning that converges within the run
($\lambda=10^{-2}$ to $10^{-4}$) the forces settle into a regular alternation in
every mode, and stance ends on the force threshold (no stance bout at the
timeout over the last 20\,s). The alternation is clearest in slow and medium walk
(corr$(F_E,F_F)\approx-0.65$ to $-0.75$ over the last 20\,s), somewhat weaker
in toe stepping ($-0.55$ to $-0.71$), and weakest in fast walk ($-0.47$ to
$-0.57$), where the short stance leaves the extensor little time to rise and
fall. Under partial unloading the gait visibly sharpens as the weight grows
(toe stepping at $\lambda=10^{-3}$: $r=-0.31$ in the first window, $-0.65$ in
the last). Extensor force peaks at $\approx4$--$7$\,a.u.: stance now ends when
fatigue has pulled force below the threshold, before it saturates.

At the two slowest rates the gait does not form. At $\lambda=10^{-5}$ the loop
still switches on the force threshold, but with the cutaneous synapse still
weak the alternation degrades (corr$(F_E,F_F)$ $-0.18$ to $-0.49$ in the walking
modes), and in toe stepping stance already runs to the timeout. At
$\lambda=10^{-6}$ the synapses stay at their initial values and, in every mode,
most stance bouts run to the timeout ($41$--$95\%$). The alternation that remains
there can look regular --- corr$(F_E,F_F)$ reaches $-0.79$ in slow walk --- but
it is imposed by the timeouts acting as a clock, not produced by the sensory
loop.

\begin{figure}[t]
  \centering
  \includegraphics[width=0.9\textwidth]{fig_force_stages_lam1em2.png}
  \caption{\textbf{Force at three learning stages, four locomotion modes,
    $\lambda=10^{-2}$ (fastest rate).} Rows: the four locomotion modes; columns:
    beginning / middle / end of the $120$\,s run. Solid = Force-E, dashed =
    Force-F (a.u.)\ versus time (s), leg L; panels annotated with the mean
    CUT$\to$RG-E weight and in-window corr$(F_E,F_F)$. Panels (a)--(l).}
  \label{fig:force_l2}
\end{figure}

\begin{figure}[t]
  \centering
  \includegraphics[width=0.9\textwidth]{fig_force_stages_lam1em3.png}
  \caption{\textbf{Force at three learning stages, four locomotion modes,
    $\lambda=10^{-3}$.} Same layout as Figure~\ref{fig:force_l2}.}
  \label{fig:force_l3}
\end{figure}

\begin{figure}[t]
  \centering
  \includegraphics[width=0.9\textwidth]{fig_force_stages_lam1em4.png}
  \caption{\textbf{Force at three learning stages, four locomotion modes,
    $\lambda=10^{-4}$.} Same layout as Figure~\ref{fig:force_l2}. At this
    intermediate rate the three stages are most distinct, spanning the
    convergence of the weights.}
  \label{fig:force_l4}
\end{figure}

\begin{figure}[t]
  \centering
  \includegraphics[width=0.9\textwidth]{fig_force_stages_lam1em5.png}
  \caption{\textbf{Force at three learning stages, four locomotion modes,
    $\lambda=10^{-5}$.} Same layout as Figure~\ref{fig:force_l2}. The
    cutaneous weight barely develops within the run and the alternation stays
    weak; in toe stepping stance runs to the timeout.}
  \label{fig:force_l5}
\end{figure}

\begin{figure}[t]
  \centering
  \includegraphics[width=0.9\textwidth]{fig_force_stages_lam1em6.png}
  \caption{\textbf{Force at three learning stages, four locomotion modes,
    $\lambda=10^{-6}$ (slowest rate).} Same layout as Figure~\ref{fig:force_l2}.
    The weights stay at their initial values; most stance bouts end at the
    timeout, so the regular pattern here is set by the timeouts rather than by
    the sensory loop.}
  \label{fig:force_l6}
\end{figure}

% =============================================================
\subsection{Network-level activity across the locomotion modes}\label{sec:results:network}

The muscle force is only the read-out of a whole network, so we next ask whether
the alternation is a property of the entire circuit.
Figures~\ref{fig:net_l2}--\ref{fig:net_l6} show the firing rates of all sixteen
recorded neuron pools --- for both legs, the rhythm generators (RG-E, RG-F), the
reciprocal-inhibition interneurons (InE, InF), the proprioceptive interneurons
(IaInt-E, IaInt-F) and the motor pools (M-E, M-F) --- over the last 4\,s of each
run, one figure per learning rate. At the converging rates every extensor-side
pool fires in alternation with its flexor-side partner in all four modes, the
two legs stay in antiphase, and the rhythm generators alternate strongly
(corr(RG-E,RG-F) between $-0.87$ and $-0.98$) --- so the alternation is built
into the core circuit, not imposed only at the muscle. The cycle shortens from
slow to fast walk while the pattern keeps its shape. Under toe stepping the
extensor-side pools (RG-E, InE, M-E) and the Ia interneurons fire at roughly
half the rate of full weight-bearing, the network-level signature of the halved
sensory feedback, while the flexor side is unchanged.

\begin{figure}[t]
  \centering
  \includegraphics[width=0.95\textwidth]{fig_network_matrix_lam1em2.png}
  \caption{\textbf{Full-circuit population activity across the four locomotion
    modes, $\lambda=10^{-2}$ (fastest rate).} Rows: the sixteen recorded
    populations (both legs; right-leg rows tinted); columns (a)--(d): slow /
    medium (plantar, baseline) / fast walk and toe stepping. Population rate
    (spikes$\cdot$neuron$^{-1}\cdot$s$^{-1}$) versus time (ms), last 4\,s of
    the run.}
  \label{fig:net_l2}
\end{figure}

\begin{figure}[t]
  \centering
  \includegraphics[width=0.95\textwidth]{fig_network_matrix_lam1em3.png}
  \caption{\textbf{Full-circuit population activity, $\lambda=10^{-3}$.} Same
    layout as Figure~\ref{fig:net_l2}.}
  \label{fig:net_l3}
\end{figure}

\begin{figure}[t]
  \centering
  \includegraphics[width=0.95\textwidth]{fig_network_matrix_lam1em4.png}
  \caption{\textbf{Full-circuit population activity, $\lambda=10^{-4}$.} Same
    layout as Figure~\ref{fig:net_l2}.}
  \label{fig:net_l4}
\end{figure}

\begin{figure}[t]
  \centering
  \includegraphics[width=0.95\textwidth]{fig_network_matrix_lam1em5.png}
  \caption{\textbf{Full-circuit population activity, $\lambda=10^{-5}$.} Same
    layout as Figure~\ref{fig:net_l2}. With the cutaneous weight still
    developing, left--right alternation is weaker, most visibly in fast walk
    and toe stepping.}
  \label{fig:net_l5}
\end{figure}

\begin{figure}[t]
  \centering
  \includegraphics[width=0.95\textwidth]{fig_network_matrix_lam1em6.png}
  \caption{\textbf{Full-circuit population activity, $\lambda=10^{-6}$
    (slowest rate).} Same layout as Figure~\ref{fig:net_l2}; stance timing is
    set mostly by the timeouts at this rate.}
  \label{fig:net_l6}
\end{figure}

% =============================================================
\subsection{Comparison across the 4 locomotion modes $\times$ 5 STDP rates}\label{sec:results:comparison}

To compare all twenty conditions at a glance, Figures~\ref{fig:mode_lambda_heat}
and~\ref{fig:mode_lambda_trend} and Table~\ref{tab:mode_lambda} reduce each
condition to six numbers measured over the last $20$\,s: extensor/flexor
alternation of the forces [corr$(F_E,F_F)$] and of the rhythm generators
[corr(RG-E,RG-F)], left--right alternation [corr$(F_{E,L},F_{E,R})$], the
fraction of stance bouts ended by the timeout, peak extensor force, and the
learned CUT$\to$RG-E weight. Three patterns stand out.

\textbf{(i) Learning is what turns the loop on.} From $\lambda=10^{-2}$ to
$10^{-4}$ stance essentially always ends on the force threshold (no stance bout
at the timeout over the last 20\,s; at most $7\%$ on one leg from $t=30$\,s),
and the legs alternate (corr$(F_{E,L},F_{E,R})$ $-0.55$ to $-0.81$, except fast
walk at $10^{-2}$, $-0.32$). At $10^{-6}$, where the synapses do not learn, the timeouts take over
in every mode ($41$--$95\%$ of stance bouts), and at $10^{-5}$ in toe stepping
($100\%$). A closed, force-triggered gait therefore requires a learned
cutaneous synapse; without one the network falls back on its safety timer.

\textbf{(ii) The worst gait is at a partially learned weight.} At
$\lambda=10^{-5}$, with CUT$\to$RG-E at $\approx7$--$14$\,pA, the loop is
closed in the walking modes but poorly coordinated: extensor/flexor alternation
is the weakest of the grid (corr$(F_E,F_F)$ $-0.18$ to $-0.49$) and left--right
alternation drops in medium and fast walk and in toe stepping. A partially
developed cutaneous drive is strong enough to take over the timing of stance but
not strong or well-timed enough to organise it; with no learning at all the
timer is more regular, and with converged learning the sensory loop is.

\textbf{(iii) Speed and unloading.} Among the converged conditions, alternation
is best in slow and medium walk and falls in fast walk, where stance lasts only
$\approx190$\,ms. Partial unloading lowers extensor/flexor alternation
moderately and slows learning, but toe stepping still produces a genuine,
alternating gait once the cutaneous weight has converged. The rhythm generators
alternate strongly in every condition (corr(RG-E,RG-F) $\le-0.82$), including
those where the muscle forces do not: the counter-phase is lost at the muscle
level before it is lost at the rhythm-generator level.

\begin{figure}[t]
  \centering
  \includegraphics[width=0.95\textwidth]{fig_mode_lambda_heatmaps.png}
  \caption{\textbf{Metric heatmaps across the $4\times5$ grid} (rows: locomotion
    modes; columns: STDP rate $\lambda$). corr$(F_E,F_F)$, corr(RG-E,RG-F) and
    corr$(F_{E,L},F_{E,R})$ (diverging scale), fraction of stance bouts ended by
    the timeout (worst leg), peak Force-E (a.u.), and learned CUT$\to$RG-E weight
    (pA). Last 20\,s of each run; leg L unless stated.}
  \label{fig:mode_lambda_heat}
\end{figure}

\begin{figure}[t]
  \centering
  \includegraphics[width=0.9\textwidth]{fig_mode_lambda_trends.png}
  \caption{\textbf{Metric trends across the locomotion modes, one line per STDP
    rate $\lambda$.} Same six metrics as Figure~\ref{fig:mode_lambda_heat},
    plotted against the locomotion mode (slow walk $\to$ toe stepping).}
  \label{fig:mode_lambda_trend}
\end{figure}

\input{mode_lambda_table}
